\documentclass[11pt,a4paper]{article}

\usepackage[utf8]{inputenc}
\usepackage[T1]{fontenc}
\usepackage[spanish]{babel}
\usepackage[a4paper,margin=2.5cm]{geometry}
\usepackage{booktabs}
\usepackage{array}
\usepackage{xcolor}
\usepackage{listings}
\usepackage[hidelinks]{hyperref}
\usepackage{enumitem}

\lstset{
  basicstyle=\ttfamily\small,
  frame=single, breaklines=true, showstringspaces=false,
  columns=fullflexible
}

\title{Bitácora de sesión --- mammoviz-tda\\\large Resumen de trabajo del 17 de julio de 2026}
\author{Sesión asistida con Claude Code}
\date{2026-07-17}

\begin{document}
\maketitle

\section*{Punto de partida}
Al empezar la sesión, el repositorio \texttt{mammoviz-tda} (TFG sobre Análisis
Topológico de Datos aplicado a la clasificación de mamografías CBIS-DDSM)
tenía:
\begin{itemize}[nosep]
  \item Memoria LaTeX con 8 capítulos, de los cuales los capítulos 7
        (Resultados) y 8 (Conclusiones) eran plantillas con marcadores de
        posición (\textit{``--''}, \textit{``Redactar a partir de los
        resultados obtenidos''}).
  \item \textit{Pipeline} en \texttt{src/} (\texttt{data.py},
        \texttt{preprocessing.py}, \texttt{topology.py}, \texttt{models.py},
        \texttt{evaluate.py}) sin ningún test automatizado.
  \item Un notebook de demostración (\texttt{01\_pipeline\_demo.ipynb}) que
        funcionaba, pero caía a un descriptor de respaldo (curva de Betti-0)
        porque \texttt{giotto-tda} no tiene binarios para Python 3.14 (versión
        instalada en esta máquina).
  \item Sin datos reales de CBIS-DDSM descargados (solo el manifiesto de TCIA).
  \item Dos figuras de la memoria eran cuadros de texto de relleno
        (\textit{``sustituir por la figura definitiva''}), no imágenes reales.
\end{itemize}

\section*{Trabajo realizado}

\subsection*{1. Capítulos 7 y 8 de la memoria}
Se sustituyeron los marcadores de posición por contenido real, basado en la
ejecución efectiva del \textit{pipeline} (no en resultados inventados):
diseño experimental (E1--E4), configuración, estado explícito de los datos
pendientes de descargar, y una tabla de cumplimiento de objetivos (O1--O6)
honesta sobre qué está implementado frente a qué falta validar con datos
reales.

\subsection*{2. Tests unitarios (\texttt{pytest})}
Se creó el directorio \texttt{tests/} con suite completa para
\texttt{src/}: \texttt{test\_data.py}, \texttt{test\_preprocessing.py},
\texttt{test\_models.py}, \texttt{test\_evaluate.py} y tests de topología
divididos en \texttt{test\_topology.py} (genérico),
\texttt{test\_topology\_gudhi.py} y \texttt{test\_topology\_giotto.py}
(específicos de cada backend, con \texttt{pytest.importorskip}).
Se detectó y corrigió un error de diseño propio: un \texttt{importorskip} a
mitad de fichero salta el \textbf{módulo entero}, no solo lo que viene
después --- de ahí la separación en varios ficheros.
Resultado final: \textbf{43 tests pasan, 2 se saltan} (los que requieren
TensorFlow, no instalado en este entorno).

\subsection*{3. Backend de homología persistente portable (gudhi)}
\texttt{giotto-tda} solo publica binarios hasta Python 3.12. En vez de
instalar una versión antigua de Python en paralelo, se identificó que
\texttt{gudhi} + \texttt{persim} sí tienen binarios para Python 3.14 y
calculan la misma homología persistente cúbica. Se reescribió
\texttt{src/topology.py} para elegir automáticamente el mejor backend
disponible (\texttt{giotto-tda} $\to$ \texttt{gudhi} $\to$ curva de Betti-0),
con una función \texttt{backend()} para consultarlo, y una implementación
propia de imagen de persistencia (núcleo gaussiano ponderado por
persistencia) para cuando el backend activo es \texttt{gudhi}.

\subsection*{4. Notebook y resultados actualizados}
Se actualizó \texttt{01\_pipeline\_demo.ipynb} para usar \texttt{backend()} y
reportar correctamente qué motor topológico está activo, y se volvió a
ejecutar de extremo a extremo. Al pasar del descriptor de respaldo (Betti-0,
solo H0) a la imagen de persistencia completa (H0+H1, backend \texttt{gudhi}),
los resultados sobre el conjunto sintético mejoraron notablemente:

\begin{table}[h]
\centering
\begin{tabular}{@{}lccc@{}}
\toprule
& \textbf{Antes (Betti-0, solo H0)} & \textbf{Después (gudhi, H0+H1)} \\
\midrule
Sensibilidad (Random Forest) & 0.60 & 0.80 \\
AUC (SVM / Random Forest)    & 0.74 / 0.78 & 0.844 / 0.833 \\
Dimensión del vector         & 24 & 288 \\
\bottomrule
\end{tabular}
\caption{Comparación de resultados sobre el conjunto sintético (120 imágenes),
antes y después de disponer de un backend TDA real.}
\end{table}

\subsection*{5. Figuras reales en la memoria}
Se sustituyeron los dos últimos cuadros de texto de relleno:
\begin{itemize}[nosep]
  \item El diagrama del \textit{pipeline} (Cap.\ 5) ahora es un diagrama TikZ
        real (código, no imagen externa), corregido tras detectar que se
        salía del margen a tamaño natural (\texttt{resizebox}).
  \item Tres figuras extraídas de la ejecución real del notebook (Cap.\ 7):
        ejemplos sintéticos benigno/maligno, matriz de confusión del mejor
        modelo, y diagramas de persistencia H0+H1 de un caso de cada clase.
\end{itemize}

\subsection*{6. Apéndice ampliado}
Se rellenó el apéndice (antes vacío/genérico) con una tabla real de los
campos del manifiesto de TCIA (verificada contra el \texttt{.xlsx} presente en
el repositorio), la descripción de los campos esperados en los CSV de casos
(aún por descargar), y dos fragmentos de código relevantes (detección de
backend, imagen de persistencia manual). De paso se detectaron y corrigieron
dos tablas con columnas sin ajuste de texto que desbordaban la página (bug
preexistente en el capítulo 8, agravado por la tabla nueva del apéndice).

\subsection*{Verificación}
Cada cambio se verificó activamente, no solo se asumió correcto:
\begin{itemize}[nosep]
  \item La memoria se recompiló (\texttt{latexmk + biber}) tras cada tanda de
        cambios, comprobando 0 errores y 0 referencias sin resolver.
  \item Las páginas con figuras y tablas nuevas se renderizaron a PNG
        (\texttt{pdftoppm}) y se inspeccionaron visualmente para detectar
        desbordamientos de márgen, no solo confiar en que ``compila''.
  \item La suite de tests se ejecutó de principio a fin varias veces
        (43 passed, 2 skipped, estable).
\end{itemize}

\section*{Estado del repositorio al cierre}
\textbf{Nada de esto se ha commiteado} (a petición explícita). Ficheros
modificados o nuevos en el árbol de trabajo:
\begin{lstlisting}
 M .gitignore
 M README.md
 M notebooks/01_pipeline_demo.ipynb
 M requirements.txt
 M src/topology.py
 M tfg/capitulos/05-metodologia.tex
 M tfg/capitulos/06-implementacion.tex
 M tfg/capitulos/07-resultados.tex
 M tfg/capitulos/08-conclusiones.tex
 M tfg/main.pdf
 M tfg/main.tex
?? conftest.py
?? pytest.ini
?? tests/
?? tfg/figuras/diagramas_persistencia_sintetico.png
?? tfg/figuras/ejemplos_sinteticos.png
?? tfg/figuras/matriz_confusion_sintetico.png
\end{lstlisting}

\section*{Pendiente (fuera del alcance de esta sesión)}
Todo lo pendiente depende de un mismo bloqueo externo: la descarga real de
CBIS-DDSM (CSV de casos + imágenes DICOM/máscaras vía NBIA Data Retriever,
del orden de decenas de GB), que requiere cuenta en TCIA y no es automatizable
desde este entorno. Una vez descargados los datos:
\begin{enumerate}[nosep]
  \item Ejecutar \texttt{src/topology.py} (ya operativo con \texttt{gudhi})
        sobre las imágenes reales y rellenar la Tabla 7.2.
  \item Entrenar la CNN de referencia (\texttt{src/models.py:build\_cnn}) ---
        requiere además instalar TensorFlow.
  \item Ejecutar el experimento de fusión TDA + CNN.
  \item Generar las visualizaciones finales (Grad-CAM vs.\ diagramas de
        persistencia) sobre casos reales.
\end{enumerate}

\end{document}
